%! Author = lili (Transkription)
%! Date = 3/7/26

\untit{Vorlesung 3. Juli 2026}

\section{Punktschätzer für den Erwartungswert: Rückblick}\label{sec:punktschaetzer-rueckblick-juli}

\txc{(2026)} Grundsätzlich ist $\bar X = \frac1n\sum_{i=1}^n X_i$ ein \gls{Schätzer}, der leicht zu berechnen ist.
Was schätzt $\bar X$? Grundsätzlich ist $\bar X$ ein erwartungstreuer \gls{Schätzer} für den EW.
Wir haben in den letzten Vorlesungen sowohl Beispiele gesehen, bei denen der \gls{param}, den wir eigentlich
schätzen wollen, gleich dem EW ist, als auch Beispiele, bei denen der \gls{param} aus dem EW ausgerechnet werden
kann.

\begin{bsp}
    \txc{(2026)} \\
    Wenn die Stichprobenvariablen $X_1,\dots,X_n$ Poisson-verteilt sind mit \gls{param} $\lambda>0$, so gilt
    \[
        \lambda = \E{X_1} \cdot
    \]
    Wenn die Stichprobenvariablen $X_1,\dots,X_n$ exponentialverteilt sind mit \gls{param} $\alpha>0$, gilt dagegen
    \[
        \alpha = \frac{1}{\E{X_1}} \cdot
    \]
\end{bsp}

Es ist nicht immer möglich, den \gls{param} einer \gls{verteilung} aus dem EW zu berechnen.

\begin{bsp}
    \txc{(2026)} \\
    Seien die Stichprobenvariablen gleichverteilt auf $\{-\theta,\dots,\theta\}$, $\theta\in\nz$, $\theta>0$.
    Dann
    gilt stets
    \[
        \E{X_i} = 0 \cdot
    \]
\end{bsp}
\komm{
    Hier könnten wir statt des EW die Varianz schätzen (die Stichprobenvariablen sind ja identisch verteilt).
    Das führt zur Frage: Wie schätzen wir Varianzen?
}

\karkar{bem_poisson_parameter_aus_ew}
\karkar{bem_exponential_parameter_aus_ew}
\karkar{bem_parameter_nicht_immer_aus_ew_gewinnbar}
\karkar{bem_alternative_wenn_ew_nicht_ausreicht}

\section{Wiederholung: Schätzer für die Varianz}\label{sec:varianzschaetzer-wiederholung}

\txc{(2026)}
\paragraph{Notation.}
Stichprobe: $\bar x = \dfrac1n\sum_{i=1}^n x_i$. \\
Stichprobenvariable: $\bar X = \dfrac1n\sum_{i=1}^n X_i$ ist erwartungstreu.

Wir dürfen annehmen, dass $\E{X_i}=0$ gilt \txo{(o.\,B.\,d.\,A., siehe letzte Vorlesung: eine Nullergänzung zeigt,
    dass $X_i - \bar X$ nicht vom tatsächlichen Erwartungswert der $X_i$ abhängt)}.

\begin{er}
    \txc{(2026)} \\
    Der naive Schätzer $s^2(x) = \dfrac1n\sum_{i=1}^n (x_i-\bar x)^2$ ist \emph{nicht} erwartungstreu:
    \[
        \E{s^2(X)} = \left(1-\frac1n\right) \Var(X_1) \cdot
    \]
    Der korrigierte Schätzer
    \[
        \widehat\sigma^2(x) = \frac{1}{n-1}\sum_{i=1}^n (x_i-\bar x)^2
    \]
    ist erwartungstreu ($s^2$ mit dem Faktor $\frac{n}{n-1}$ multipliziert).
\end{er}

Wenn der Erwartungswert bekannt ist, können wir den naiven Ansatz im Prinzip verwenden.
Wir müssen $\bar x$ durch
den EW ersetzen.

\begin{satz}
    \txc{(2026)} \\
    \textbf{Schätzer für die Varianz bei bekanntem EW $\wfk$:}
    \[
        \frac1n \sum_{i=1}^n (x_i-\wfk)^2
    \]
    ist tatsächlich erwartungstreu.
\end{satz}

\begin{bsp}
    \txc{(2026)} \\
    \textbf{Illustration mit simulierten Daten (Mathematica).}
\end{bsp}
\komm{
    In der Vorlesung wurde dies an einem Beispiel mit $n=100$ unabhängig identisch verteilten Stichproben aus
    einer Exponentialverteilung mit \gls{param} $\alpha=2$ illustriert (zwei unabhängige Durchläufe mit
    unterschiedlichem Seed). Berechnet wurden jeweils: der Schätzer für den EW ($\bar x$), der Schätzer für die
    Varianz bei unbekanntem EW ($\widehat\sigma^2$, Nenner $n-1$), sowie -- im zweiten Durchlauf -- zusätzlich der
    Schätzer für die Varianz bei bekanntem EW $\wfk=1/\alpha$ (Nenner $n$) zum Vergleich mit $\bar x^{\,2}$. In
    beiden Durchläufen lagen die Schätzwerte nahe, aber nicht exakt bei den wahren Werten $1/\alpha=0.5$ bzw.\
    $1/\alpha^2=0.25$ -- die genauen Zahlenwerte aus der Mitschrift ließen sich aus der Quelle nicht zuverlässig
    rekonstruieren.
}

\karkar{satz_varianzschaetzer_bekannter_ew}

\section{Die Normalverteilung}\label{sec:normalverteilung-uebergang}

\txc{(2026)} Um zu verstehen, wie ein Schätzer für eine Varianz sinnvoll sein kann, wollen wir uns als Beispiel
die Normalverteilung ansehen, die sich von unseren bisherigen Beispielen dadurch unterscheidet, dass sie durch
zwei \gls{param} beschrieben wird.

\begin{defi}
    \textbf{Dichte der Normalverteilung} \txc{(2026)} \\
    Für $\wfk\in\rz$, $\sigma^2>0$ ist die Dichte der Normalverteilung gegeben durch
    \[
        \varphi_{\wfk,\sigma^2}(x) = \frac{1}{\sqrt{2\pi\sigma^2}} \exp\left( -\frac{(x-\wfk)^2}{2\sigma^2} \right)
        \qquad \forall\, x\in\rz \cdot
    \]
    Für $\wfk=0,\ \sigma^2=1$ ist dies die \textbf{Dichte der Standardnormalverteilung}
    \[
        \varphi(x) = \frac{1}{\sqrt{2\pi}} e^{-x^2/2} \qquad \forall\, x\in\rz \cdot
    \]
\end{defi}

\begin{satz}
    \textbf{Fakten} \txc{(2026)} \\
    Ist $X$ eine \gls{zv} mit Dichte $\varphi_{\wfk,\sigma^2}$, so gelten
    \[
        \E{X} = \wfk \qquad \text{und} \qquad \Var(X) = \sigma^2 \cdot
    \]
    Daher sprechen wir bei $\varphi_{\wfk,\sigma^2}$ von der Dichte einer Normalverteilung mit EW $\wfk$ und
    Varianz $\sigma^2$.
\end{satz}

\begin{bsp}
    \txc{(2026)} \\
    \textbf{Illustration mit simulierten Daten (Mathematica).}
\end{bsp}
\komm{
    Auch hier zeigte die Vorlesung ein Beispiel mit $n=100$ i.i.d.\ Stichproben, diesmal aus $N(\wfk=20,\sigma^2=1)$.
    Berechnet wurden der Schätzer für den EW, der Schätzer für die Varianz bei unbekanntem EW ($\widehat\sigma^2$)
    sowie der Schätzer bei bekanntem EW $\wfk=20$ zum Vergleich. Die konkreten Zahlenwerte waren in der
    Quelle nicht zuverlässig lesbar.
}

\karkar{bem_warum_normalverteilung_eingefuehrt}
\karkar{def_dichte_normalverteilung_allgemein}
\karkar{bem_fakten_ew_var_normalverteilung}

\section{Bereichsschätzer: Konfidenzintervalle}\label{sec:bereichsschaetzer-konfidenzintervalle}

\txc{(2026)} Die Beispiele zeigen, dass der geschätzte \gls{param} in der Regel nicht der wahre \gls{param} ist.

\paragraph{Idee.} Wir suchen ein Intervall, so dass wir mit vorgegebener \gls{ws} sagen können, dass der \gls{param}
in diesem Intervall liegt.

\subsection*{Erste Überlegungen}

Sei $Z\sim N(0,1)$ standardnormalverteilt.
Suche $c$ derart, dass
\[
    \alpha \geq \pe(|Z|>c)
\]
für ein gegebenes $\alpha\in(0,1)$.
\[
    \alpha \geq \pe(|Z|>c) \Leftrightarrow \alpha \geq \Phi(-c) + (1-\Phi(c)) = 2(1-\Phi(c)) \Leftrightarrow \Phi(c) \geq
    1 - \frac\alpha2 \cdot
\]
Als \gls{verteilung}sfunktion ist $c\mapsto\Phi(c)$ monoton wachsend.
Wir können daher das kleinste $c$ mit
$\Phi(c)\geq 1-\frac\alpha2$ bestimmen -- verwende eine Tabelle, um dieses $c$ zu finden.

\begin{defi}
    \txc{(2026)} \\
    Das kleinste $c$ mit $\Phi(c) \geq 1-\frac\alpha2$ heißt das \textbf{$(1-\alpha/2)$-Quantil der
    Standardnormalverteilung}, geschrieben $c = z_{1-\alpha/2}$.
\end{defi}

\begin{bsp}
    \txc{(2026)} \\
    $\alpha = 0.05$: $\Leftrightarrow \Phi(c) \geq 0.975$ (Tabelle) $\approx c = z_{0.975} \approx 1.96$.
\end{bsp}

\subsection*{Vorgehen für die einfache Normalverteilung}

Sei $X\sim N(\wfk,\sigma^2)$ (mit bekanntem $\sigma^2$). Dann können wir $X$ mit Hilfe von $Z\sim N(0,1)$
darstellen als
\[
    X = \wfk + \sigma Z \cdot
\]
Somit gilt
\[
    \pe(\wfk \in [X-\sigma c,\ X+\sigma c]) = \pe(|Z|\leq c) \geq 1-\alpha \cdot
\]
Insbesondere:
\[
    \pe(\wfk \in [X-1.96\,\sigma,\ X+1.96\,\sigma]) \geq 0.95 \cdot
\]

\subsection*{Allgemein}

Seien nun $X_1,\dots,X_n$ u.i.v.\ Stichprobenvariable, $X_i \sim N(\wfk,\sigma^2)$.

\begin{satz}
    \textbf{Summe unabhängiger normalverteilter \gls{zv}} \txc{(2026)} \\
    Die Summe unabhängiger normalverteilter \gls{zv} ist wieder normalverteilt.
    Die \gls{param} addieren sich:
    \[
        \sum_{i=1}^n X_i \sim N(n\wfk,\ n\sigma^2) \cdot
    \]
    Folglich
    \[
        \bar X \sim N\!\left(\wfk,\ \frac{\sigma^2}{n}\right) \cdot
    \]
\end{satz}

Damit gilt für das \textbf{Konfidenzintervall für $\wfk$ bei bekanntem $\sigma^2$}:
\begin{equation}
    \label{eq:ki-bekannte-varianz}
    \pe\!\left( \wfk \in \left[ \bar X - \sigma\sqrt{\tfrac1n}\, c,\ \ \bar X + \sigma\sqrt{\tfrac1n}\, c \right]
    \right) \geq 1-\alpha \cdot
\end{equation}
Mit wachsender Stichprobengröße $n$ wird das Intervall kleiner.

\subsection*{Verwendung für Versuchsplanung}

Wie viele Werte sollte meine Stichprobe haben, damit das Intervall eine vorgegebene Größe nicht überschreitet?
\komm{
    D.\,h.\ konkret: Ist eine maximale Halbbreite $d$ vorgegeben, so löst man $\sigma\sqrt{1/n}\,c \leq d$ nach $n$
    auf und erhält $n \geq \left(\dfrac{\sigma c}{d}\right)^2$.
}

\begin{defi}
    \textbf{Konfidenzintervall} \label{defi:konfidenzintervall} \txc{(2026)} \\
    Sei $\theta$ der wahre \gls{param} der \gls{verteilung}.
    Seien $U,V$ Funktionen von $X_1,\dots,X_n$ (also
    Funktionen von $\eraum$ nach $\rz$). Wenn
    \[
        \pe_{\theta}\big( \theta \in [U,V] \big) \geq 1-\alpha \qquad \forall\, \theta
    \]
    gilt, dann heißt $[U,V]$ ein \textbf{Konfidenzintervall} (oder \textbf{Vertrauensintervall}) für den
    unbekannten \gls{param} $\theta$ zum \textbf{Konfidenzniveau} $1-\alpha$.
\end{defi}

\subsection*{Konfidenzintervalle in Standardsituationen}

\paragraph{Stichprobe aus einer Normalverteilung mit bekannter Varianz $\sigma^2$.}
Das Konfidenzintervall für den Erwartungswert ist gegeben durch
\[
    \left[\ \bar X - \sigma\sqrt{\tfrac1n}\, c,\ \ \bar X + \sigma\sqrt{\tfrac1n}\, c\ \right]
\]
mit $c$ dem Quantil der Standardnormalverteilung.

\paragraph{Stichprobe aus einer Normalverteilung mit unbekannter Varianz.}
In diesem Fall müssen wir auch die Varianz schätzen.
Daher gilt: Das Konfidenzintervall für den Erwartungswert ist
gegeben durch
\[
    \left[\ \bar X - \widehat\sigma\sqrt{\tfrac1n}\, t_{n-1,1-\alpha/2},\ \ \bar X + \widehat\sigma\sqrt{\tfrac1n}\,
        t_{n-1,1-\alpha/2}\ \right]
\]
Dabei ist $\widehat\sigma^2 = \dfrac{1}{n-1}\displaystyle\sum_{i=1}^n (x_i-\bar x)^2$ der Schätzer der Varianz, und
$t_{n-1,1-\alpha/2}$ das $(1-\alpha/2)$-Quantil der sogenannten \textbf{Studentschen $t$-\gls{verteilung}} mit $n-1$
Freiheitsgraden.

\begin{defi}
    \textbf{Studentsche $t$-\gls{verteilung}} \label{defi:studentsche-t-verteilung} \txc{(2026)} \\
    Die Studentsche $t$-\gls{verteilung} mit $n-1$ Freiheitsgraden ist die \gls{verteilung} der \gls{zv}
    \[
        T = \frac{\bar X - \wfk}{\widehat\sigma / \sqrt n} \cdot
    \]
    Diese \gls{zv} hat eine Dichte, und ihre \gls{verteilung}sfunktion ist gegeben durch das Integral über diese Dichte.
    Die Normierungskonstante der Dichte kann durch die Gamma-Funktion ausgedrückt werden.
\end{defi}
\komm{
    Die genaue Dichteformel über die Gamma-Funktion war in der vorliegenden Quelle nicht zuverlässig
    rekonstruierbar und wurde deshalb bewusst weggelassen -- bei Bedarf bitte im Original-Skript nachschlagen.
}

Wie bei der Normalverteilung sind wir für die Werte der \gls{verteilung}sfunktion auf eine Tabelle angewiesen.
\komm{
    Die Vorlesung zeigte hierzu Dichte-Plots der $t$-\gls{verteilung} für verschiedene Freiheitsgrade (u.\,a.\ 2, 3, 5,
    10, 30) im Vergleich zur Standardnormalverteilung: Für $n\to\infty$ nähert sich die $t$-\gls{verteilung} erkennbar
    der Standardnormalverteilung an.
}

\subsection*{Konfidenzintervall für eine Bernoulli-verteilte Stichprobe}

Wenn der Erwartungswert zugleich die Erfolgswahrscheinlichkeit geschätzt werden soll, dann kennen wir die Varianz
nicht direkt.
Wir nutzen $\bar X$ als Schätzer für den EW und $\widehat\sigma^2$ als Schätzer für die Varianz.

\begin{er}
    \txc{(2026), Hausaufgabe} \\
    Es gilt näherungsweise
    \[
        \widehat\sigma^2 \approx \bar X (1-\bar X)
    \]
    für große $n$.
\end{er}

Mit der Normalapproximation ergibt sich das Konfidenzintervall für die Erfolgswahrscheinlichkeit wie im Fall der
unbekannten Varianz.
Für große $n$ kann $t_{n-1,1-\alpha/2}$ durch $z_{1-\alpha/2}$ approximiert werden (die
$t$-\gls{verteilung} nähert sich für große $n$ der Standardnormalverteilung an; siehe Bilder oben).

Damit finden wir:
\begin{equation}
    \label{eq:ki-bernoulli}
    \left[\ \bar X - \sqrt{\tfrac{\widehat\sigma^2}{n}}\, z_{1-\alpha/2},\ \ \bar X + \sqrt{\tfrac{\widehat\sigma^2}
    {n}}\, z_{1-\alpha/2}\ \right] \cdot
\end{equation}

\komm{
    In den Hausaufgaben werden derartige Konfidenzintervalle (approximativ via $t$-\gls{verteilung} bzw.\ via
    Normalverteilung, sowie exakt) miteinander verglichen.
}

\karkar{bem_motivation_bereichsschaetzer}
\karkar{def_quantil_standardnormalverteilung_ki}
\karkar{bem_darstellung_normalverteilte_zv_via_standardnormal}
\karkar{satz_summe_unabhaengiger_normalverteilter_zv_ki}
\karkar{formel_ki_bekannte_varianz}

\karkar{bem_intervallbreite_und_stichprobengroesse}
\karkar{formel_versuchsplanung_stichprobengroesse}
\karkar{def_konfidenzintervall_formal}
\karkar{bem_intuition_konfidenzniveau}
\karkar{formel_ki_unbekannte_varianz}

\karkar{bem_warum_t_verteilung_statt_normal}
\karkar{def_studentsche_t_verteilung_ki}
\karkar{bem_eigenschaften_t_verteilung_ki}
\karkar{bem_bernoulli_varianzschaetzer_approximation}
\karkar{formel_ki_bernoulli}

\karkar{bem_warum_z_statt_t_bei_bernoulli_grosses_n}